% ENG 2440 Assignment 2 report. Build: tectonic main.tex  (figures are read from ../figures/)
\documentclass[11pt,a4paper]{article}
\usepackage[margin=2cm]{geometry}            % figures are placed at their drawn width (6.30 in), not scaled
\usepackage{newtxtext,newtxmath}             % Times, the same typeface as the figures
\usepackage{graphicx,booktabs,multirow,array,xcolor,microtype}
\usepackage[font=small,labelfont=bf,skip=5pt]{caption}
\usepackage[compact]{titlesec}
\usepackage{enumitem}
\usepackage[hidelinks]{hyperref}
\usepackage{placeins}
\usepackage{capt-of}

\graphicspath{{../figures/}}
\newcommand{\fig}[2][1]{\includegraphics[width=#1\dimexpr 6.3in\relax]{#2}}   % figure at (a fraction of) its drawn width
\newcommand{\repourl}{https://github.com/DannyZ061205/ENG2440-A2}
\setlist{nosep,leftmargin=*}
\setlength{\emergencystretch}{1.5em}
\setcounter{topnumber}{3}\setcounter{bottomnumber}{2}\setcounter{totalnumber}{5}
\renewcommand{\topfraction}{0.9}\renewcommand{\bottomfraction}{0.6}\renewcommand{\textfraction}{0.07}
\setlength{\textfloatsep}{10pt plus 3pt minus 2pt}
\setlength{\floatsep}{8pt plus 2pt minus 2pt}
\setlength{\intextsep}{12pt plus 2pt minus 2pt}
\titlespacing*{\section}{0pt}{10pt}{4pt}
\titlespacing*{\subsection}{0pt}{8pt}{3pt}
\linespread{1.025}\selectfont
\setlength{\parskip}{1pt plus 1pt}

\begin{document}

\begin{center}
{\Large\bfseries Whole-tumour segmentation in multimodal 3D brain MRI}\\[3pt]
Zichen Zhao (25011189) \quad·\quad ENG 2440 Medical Imaging \& AI in Healthcare, Assignment 2\\[1pt]
Code and executed notebook: \url{\repourl}
\end{center}

\noindent\textbf{Summary.} We built a PyTorch pipeline that segments the whole brain tumour in four-channel 3D MRI (MSD
Task01, 484 BraTS cases, supplied case-level split of 338/73/73). A lightweight 3D U-Net (5.65\,M parameters) is trained on
$64^3$ patches and applied to whole volumes with an overlapping sliding window. In a controlled comparison over three
seeds, tumour-aware patch sampling raised validation sensitivity by 0.029 but added 5.3\,mL of false-positive volume and
did not raise Dice ($-$0.006, 95\% CI $-$0.015 to 0.003), so uniform sampling was kept. On the untouched test split the final
model reaches a mean Dice of 0.885 (95\% CI 0.869 to 0.899), a median HD95 of 4.9\,mm and a median absolute volume error
of 10.1\%. Most remaining errors are boundary disagreements a voxel or two wide; the clearest failures are missed small
lesions and context errors, which no single metric captures. The reference is a research segmentation, so these numbers
measure agreement with expert outlines, not clinical suitability.

\section{Task and data}
\textbf{Task.} The input is one patient case of the Medical Segmentation Decathlon (MSD) Task01 BrainTumour data
\cite{antonelli2022,simpson2019}: four co-registered MRI volumes (FLAIR, T1-weighted, contrast-enhanced T1-weighted (T1ce)
and T2-weighted), stored as one 4-channel NIfTI array of $240\times240\times155$ voxels. The output is a binary mask on the
same grid that marks the \emph{whole tumour}. Four units must be kept apart: a \emph{voxel} is one
$1\times1\times1$\,mm$^3$ cell; a \emph{slice} is one 2D plane of voxels; a \emph{volume} is the 3D array of one channel or
of the label; and a \emph{patient case} is one patient's four volumes plus their label. Training draws 3D patches from
volumes, but every metric is computed per case on the whole volume, so the case is the unit of analysis. The intended use
is technical: a research pipeline that outlines the tumour and measures its volume.

\textbf{What the label means.} The reference comes from the BraTS 2016/2017 challenges \cite{menze2015,bakas2017}, in which
raters outlined oedema (label 1), non-enhancing tumour (2) and enhancing tumour (3). We map every non-zero label to 1
and keep the original label files unchanged. The reference is a research segmentation, not pathology-confirmed tumour
extent, and raters disagree most at the diffuse edge of the oedema. Our scores therefore measure agreement with expert
outlines; they say nothing about clinical safety or suitability for treatment planning.

\textbf{Data and geometry.} We use the 484 labelled cases with the supplied \texttt{assignment2\_split.csv} unchanged:
338 training, 73 validation and 73 test cases. Each case is one patient, so the split is by patient, and assertions confirm
that the three sets of case identifiers are disjoint and that all 968 files exist. The headers of all 484 images and labels agree: the same
$240\times240\times155$ grid of 1\,mm isotropic voxels, an identity affine with orientation code RAS, and an image affine equal
to its label affine. The four channels are stored in the order FLAIR, T1, T1ce, T2 (from \texttt{dataset.json}), as
\texttt{float32}, and the labels as \texttt{uint8} with values $\{0,1,2,3\}$. The median reference tumour volume is 95.5,
91.6 and 89.1\,mL in the three splits (Appendix Table~\ref{tab:data}, Figure~\ref{fig:split}); the test labels were opened at this point only to count
voxels.

\textbf{Index order is not anatomy.} Figure~\ref{fig:views} shows one case at one physical point, the tumour's centre of
mass, converted to millimetres with the affine and back. Panels (f)--(h) lay the planes out with the affine; panels (i) and
(j) plot raw array slices, as \texttt{imshow(array[:,:,k])} would, and show the brain rotated and on its side. An array
stored in another orientation (for example LPS) would hold the mirror image with exactly the same index layout. Only the
affine ties indices to anatomy, so it must be read before a slice is displayed, a flip is used for augmentation, or a
distance is measured in millimetres. Here the anatomy confirms the header where anatomy can: the cerebellum lies at low
axis-1 indices (posterior) and the vertex at high axis-2 indices (superior). Left and right cannot be checked in a
near-symmetric brain.

\begin{figure}[t]
\centering\fig{A_views.pdf}
\caption{\textbf{The same voxel index shows the same anatomy in all four channels, but only the affine says which way is
up or left.} Training case BRATS\_001 at the tumour's centre of mass (voxel (157, 116, 67) = (157, 116, 67)\,mm). (a)--(d)
Axial FLAIR, T1, T1ce and T2; (e) the original reference labels (three sub-regions). (f)--(h) The binary whole-tumour
reference (gold outline) in the three planes, radiological convention (patient's right on the viewer's left), letters
mark anatomical directions. (i), (j) The same data plotted by raw array index. Every panel shows a $192\times192$\,mm field.}
\label{fig:views}
\end{figure}

\section{Methods}
\textbf{Preprocessing.} Because every case shares one grid, no resampling is needed and the target grid is the native
1\,mm RAS grid. The only spatial steps are a reorientation to RAS (a no-op here) and cropping each case to the bounding box
of its brain (voxels non-zero in any channel), which keeps 36\% of the volume and discards only empty background; no case
needed padding to the 64-voxel patch size. The crop box is stored, and every prediction is written back into the full
grid before it is scored. Interpolation happens only in the rotation augmentation, with linear interpolation for MRI and
nearest neighbour for labels: on BRATS\_001, a $7^\circ$ rotation with linear interpolation turns the four label values into
19,191 distinct values such as 1.5, which are not valid classes, whereas nearest neighbour keeps exactly $\{0,1,2,3\}$.
Each channel of each case is then z-scored with the mean and SD of its non-zero voxels; the zero background, about 83\% of
the full volume, would otherwise drag the statistics towards 0. Afterwards every channel has brain mean 0.000 and SD 1.000
in all 484 cases, and the background stays exactly 0. The statistics come from the case itself, so no information passes
between cases or splits.

\textbf{Augmentation.} Training patches only are augmented, each transform imitating a real variation: a left--right flip
($p=0.5$; gliomas occur in either hemisphere and a mirrored brain is still plausible, whereas anterior--posterior or
superior--inferior flips would create impossible anatomy), an axial rotation of up to $\pm10^\circ$ ($p=0.5$; residual head
tilt), and a per-channel intensity scale of $\times$U(0.9,\,1.1) and shift of U($-$0.1,\,0.1) SD on brain voxels (scanner
gain left after z-scoring). Appendix Figure~\ref{fig:augment} shows one patch before and after.

\textbf{Patches and samplers.} A PyTorch \texttt{Dataset} returns patches of $[C,D,H,W]=[4,64,64,64]$ with masks
$[1,64,64,64]$; patch $i$ is drawn with its own random generator seeded by (seed, $i$), so the patch stream is fixed by the
seed alone. The \emph{uniform} sampler places the patch uniformly where it fits in the cropped volume. The
\emph{tumour-aware} sampler centres it on a randomly chosen reference tumour voxel with probability 0.5 and otherwise draws it
uniformly. We chose 64-voxel patches after measuring that, with 96-voxel patches, 99\% of uniform patches already contain
tumour (the brain box is only about $138\times169\times138$ voxels), which would leave the two samplers nothing to differ on;
at 64 voxels the share is 78\%.

\textbf{Model and training.} A lightweight 3D U-Net \cite{ronneberger2015,cicek2016} with 4 input channels and 1 output
channel (a whole-tumour logit) has five levels of 16, 32, 64, 128 and 256 feature channels, two $3\times3\times3$
convolutions per level with instance normalisation \cite{ulyanov2016} and leaky ReLU, strided convolutions for
downsampling and transposed convolutions for upsampling: 5,647,681 parameters. The loss is soft Dice \cite{milletari2016}
plus binary cross-entropy (BCE) over the batch; the optimiser is AdamW \cite{loshchilov2019} (peak learning rate $10^{-3}$,
weight decay $10^{-5}$, 100 warm-up iterations, then cosine decay to 0); the budget is 3,000 iterations of 4 patches
(12,000 patches, about 36 per training case). Each sampler is trained with seeds 0, 1 and 2, and for each seed the two
samplers form a controlled pair: both runs load the same initial-weights file, see the same case sequence and use
identical settings, and the notebook asserts that their saved
configurations differ only in the sampler. Every 300 iterations each model predicts all 73 validation cases and 15 fixed
training cases with the full inference pipeline, and the checkpoint with the best mean validation Dice is kept.

\textbf{Inference and postprocessing.} The cropped volume is covered by $64^3$ windows with stride 32 (50\% overlap), the
last window on each axis shifted to end at the border. Each window's sigmoid probabilities are weighted by a 3D Gaussian
($\sigma=8$ voxels) \cite{isensee2021}, summed and divided by the summed weights; only the blended probability is
thresholded at 0.5. As postprocessing we remove every 26-connected component smaller than $k$\,mL, with $k$ chosen on
validation; unlike keeping only the largest component, this keeps a second lesion of a multifocal tumour.

\textbf{Metrics.} Per test case we report Dice, intersection over union (IoU), sensitivity, and the 95th-percentile
Hausdorff distance (HD95), the larger of the two directed 95th-percentile distances between the mask surfaces in
millimetres \cite{taha2015,maierhein2024}. Volume is the voxel count times the voxel volume from the header (1\,mm$^3$, so
1,000 voxels = 1\,mL). Means come with 95\% percentile-bootstrap confidence intervals (CIs) over cases (2,000 resamples), and
paired differences with paired bootstrap CIs. Every choice (sampler, seed, checkpoint, rule, overlap) was made on validation, and the test labels were used
only for the final evaluation. An earlier execution with seed 0 alone also evaluated its selected model on the test
split; the two extra seeds had been planned before it, and no setting was changed in response to it.

\section{Results}
\textbf{Training behaviour.} Both samplers learn steadily and neither overfits within the budget (Figure~\ref{fig:curves}).
Averaged over three seeds, validation Dice passes 0.80 by iteration 1,500 with uniform sampling and by 2,100 with
tumour-aware sampling. For the uniform runs, training and validation curves stay close to the end (whole-volume loss 0.179
versus 0.195 at iteration 3,000), and validation Dice is still rising over the last 600 iterations (0.855 to 0.864) as
the learning rate decays to zero, which suggests mild underfitting. The tumour-aware runs vary much more between seeds in
mid-training (validation Dice SD 0.052 at iteration 1,800, against 0.005), which points to less stable optimisation. Their
patch loss is lower only because their patches hold more tumour, which shrinks the soft-Dice term; on whole volumes, the
comparable measure, the uniform runs have the lower loss throughout.

\begin{figure}[t]
\centering\fig{D_curves.pdf}
\caption{\textbf{Both samplers converge within 3,000 iterations without overfitting; tumour-aware runs are less stable
between seeds.} (a) Dice + BCE loss and (b) Dice on whole volumes with the full inference pipeline, every 300
iterations: validation (thick, 73 cases) and training (thin, 15 fixed training cases). (c) Loss on the training patches
themselves (moving average over 100 iterations), not comparable between samplers. Lines: mean over seeds 0--2; bands:
lowest to highest seed.}
\label{fig:curves}
\end{figure}

\textbf{Uniform versus tumour-aware sampling.} Tumour-aware sampling does what it is designed to do, and that is the
problem here (Table~\ref{tab:samplers}, Figure~\ref{fig:samplers}). It raises the share of patches that contain tumour
from 78.1\% to 89.0\% and more than doubles the tumour share of patch voxels (7.1\% to 15.7\%). The network learns this
inflated prior and, on whole volumes where tumour is again a small share of the brain, calls tumour more readily:
sensitivity rises by 0.029 (95\% CI 0.021 to 0.038) and missed volume falls by 3.2\,mL, but false-positive volume rises
by 5.3\,mL (3.3 to 8.0), precision falls by 0.038 and mean HD95 grows by 6.4\,mm, so Dice does not improve ($-$0.006,
CI $-$0.015 to 0.003). Oversampling pays off when uniform patches rarely contain the target; with whole tumours of
median 90\,mL and $64^3$ patches, most uniform patches already do, so the benefit is small and the shift in prior
dominates. Pushed further, the network would almost never see the hard negatives far from the tumour (bright FLAIR from
other causes, vessels, the brain's edge) and would produce false positives exactly there, which is why moderate rates,
such as nnU-Net's one in three \cite{isensee2021}, are the norm.

\begin{table}[t]
\centering\small
\caption{\textbf{Tumour-aware sampling trades false positives for sensitivity and does not raise Dice.} Validation, 73
cases, best checkpoint of each run. Columns 2--3: mean $\pm$ SD over seeds 0--2 of the case-mean. Difference:
tumour-aware $-$ uniform, paired over cases after averaging each case over seeds, with a 95\% percentile-bootstrap CI
(2,000 resamples). Patch statistics over all 12,000 training patches of each run.}
\label{tab:samplers}
\begin{tabular}{lccc}
\toprule
 & uniform & tumour-aware & difference [95\% CI] \\
\midrule
patches containing tumour (\%) & 78.1 $\pm$ 0.2 & 89.0 $\pm$ 0.1 & \\
tumour share of patch voxels (\%) & 7.2 $\pm$ 0.1 & 15.7 $\pm$ 0.1 & \\
\midrule
Dice $\uparrow$ & 0.865 $\pm$ 0.005 & 0.858 $\pm$ 0.005 & $-$0.006 [$-$0.015, +0.003] \\
sensitivity $\uparrow$ & 0.878 $\pm$ 0.011 & 0.907 $\pm$ 0.010 & +0.029 [+0.021, +0.038] \\
precision $\uparrow$ & 0.868 $\pm$ 0.014 & 0.831 $\pm$ 0.011 & $-$0.038 [$-$0.052, $-$0.025] \\
false-positive volume (mL) $\downarrow$ & 11.9 $\pm$ 1.7 & 17.2 $\pm$ 1.3 & +5.3 [+3.3, +8.0] \\
false-negative volume (mL) $\downarrow$ & 12.1 $\pm$ 0.9 & 8.8 $\pm$ 1.0 & $-$3.2 [$-$4.2, $-$2.3] \\
\bottomrule
\end{tabular}
\end{table}

\begin{figure}[t]
\centering\fig{E_comparison.pdf}
\caption{\textbf{Per case, tumour-aware sampling raises sensitivity but adds false-positive volume.} Validation, 73
cases, each case averaged over seeds 0--2; dots are cases (grey lines join the same case), bars are medians; $\Delta$:
mean paired difference (tumour-aware $-$ uniform) with its 95\% bootstrap CI. False-positive volume on a log scale,
values below 0.3\,mL drawn at 0.3.}
\label{fig:samplers}
\end{figure}

\textbf{Final configuration.} By the rule fixed beforehand (higher mean validation Dice over seeds, 0.865 against
0.858), the uniform sampler is selected, and within it seed 2, whose checkpoint has the best validation Dice (0.870). On
validation, removing connected components smaller than $k$ raises mean Dice steadily up to the largest value of the grid,
$k = 5$\,mL (0.870 to 0.873), and lowers mean HD95 from 14.5 to 10.8\,mm, whereas keeping only the largest component
lowers Dice and sensitivity because it deletes true second lesions. As the optimum lies at the edge of the grid, a
larger $k$ might score higher; we did not extend the grid after seeing test results.

\textbf{Test performance.} On the untouched test split the final model reaches a mean Dice of 0.885 (95\% CI 0.869 to
0.899; median 0.901) after postprocessing, with a median HD95 of 4.9\,mm and no empty predictions (Table~\ref{tab:test}).
The rule changes the prediction of 59 of the 73 cases but barely moves the summary: Dice rises by 0.0023 (CI 0.0004 to
0.0052), while HD95 changes by +0.47\,mm (CI $-$0.62 to 1.87), so its validation gain in HD95 does not carry over. Its
cost is hidden in the average: 67 of the 310 removed components overlapped the reference. Volumes agree closely over
the whole range (Appendix Figure~\ref{fig:volumes}): the mean signed error is $-$1.1\%, so there is no systematic bias,
but single cases spread from $-$33.4\% to +31.2\% (95\% limits of agreement), with a median absolute error of 10.1\%.

\begin{table}[t]
\centering\small
\caption{\textbf{The final model reaches a test Dice of 0.885; postprocessing changes the summary little.} Test split,
73 cases, final model (uniform sampler, seed 2); mean $\pm$ SD and median over cases; 95\% percentile-bootstrap CI of the
mean (2,000 resamples). HD95: 95th-percentile symmetric surface distance. Volume error: predicted $-$ reference volume.}
\label{tab:test}
\begin{tabular}{lccccc}
\toprule
 & \multicolumn{2}{c}{before postprocessing} & \multicolumn{3}{c}{after postprocessing} \\
\cmidrule(lr){2-3}\cmidrule(lr){4-6}
 & mean $\pm$ SD & median & mean $\pm$ SD & median & 95\% CI of mean \\
\midrule
Dice $\uparrow$ & 0.883 $\pm$ 0.071 & 0.901 & 0.885 $\pm$ 0.066 & 0.901 & [0.869, 0.899] \\
IoU $\uparrow$ & 0.796 $\pm$ 0.104 & 0.821 & 0.799 $\pm$ 0.099 & 0.821 & [0.776, 0.822] \\
sensitivity $\uparrow$ & 0.881 $\pm$ 0.095 & 0.898 & 0.880 $\pm$ 0.095 & 0.898 & [0.857, 0.901] \\
HD95 (mm) $\downarrow$ & 7.2 $\pm$ 7.1 & 4.4 & 7.7 $\pm$ 9.2 & 4.9 & [5.8, 10.0] \\
empty predictions & \multicolumn{2}{c}{0} & \multicolumn{2}{c}{0} & \\
\midrule
$|$volume error$|$ (mL) $\downarrow$ & & & 12.7 $\pm$ 12.6 & 8.0 & \\
$|$volume error$|$ (\%) $\downarrow$ & & & 12.9 $\pm$ 10.3 & 10.1 & \\
\bottomrule
\end{tabular}
\end{table}

\textbf{Where the errors are.} Most errors are disagreements about the boundary. Across the test cases the median share
of error voxels within 2\,mm of the reference surface is 92\%, and in the best case the errors form a one-voxel rim
around a correct tumour (Figure~\ref{fig:errors}a), where oedema fades into normal tissue and the reference itself is
least certain. Small structures are the second failure: of 480 reference components smaller than 1\,mL, the prediction
touches only 9\%, and postprocessing can turn partly detected small lesions into missed ones. In BRATS\_253 the model
touched three satellite lesions, 60 to 82\,mm from the main tumour, only with tiny fragments (0.2\,mL, half of it true
positive); the rule deleted them, Dice barely changed (0.902 to 0.903), but HD95 rose from 19.4 to 56.9\,mm. In the
median case (Figure~\ref{fig:errors}b) every false positive is attached to the tumour (23.3\,mL along its border), and
most of the 12.1\,mL it misses lies where the tumour extends across the midline. The worst case is a patch-context
failure (Figure~\ref{fig:errors}c): its coronal and sagittal planes are smeared, consistent with thick original slices,
and the prediction grows from the true lesion into the posterior fossa as one connected region, adding 53.3\,mL of false
positives while missing 25.0\,mL of the lesion. A $64^3$ patch sees only local texture; where that texture looks abnormal
for other reasons, nothing in the patch tells the network where the tumour ends.

\begin{figure}[t]
\centering\fig{G_errors.pdf}
\caption{\textbf{Errors are mostly thin boundary rims; the median case under-calls the part of the tumour that crosses the
midline, and the worst case's false positives spread into the posterior fossa.} Test cases with the highest, median and lowest Dice
(chosen by rule), final model after postprocessing, on FLAIR. Columns 1--3: planes through the centre of mass of
reference $\cup$ prediction, radiological convention, same field of view and scale (bar: 5\,cm). Column 4: all axial
slices projected onto one image (each pixel coloured by the more frequent error type along its column).}
\label{fig:errors}
\end{figure}

\textbf{One metric is not enough.} The cases rank differently on each metric. The median case's HD95 (9.8\,mm) is only 1.6\,mm
smaller than the worst case's (11.4\,mm) although its Dice is 0.30 higher, because HD95 reports the farthest 5\% of
boundary errors and both cases have errors about a centimetre out. Volume error separates them more (+6.4\% against
+33.2\%) but hides compensating mistakes: missing 25\,mL and adding 25\,mL scores 0\%. BRATS\_253 has a Dice near the
median (0.903) but an HD95 of 56.9\,mm, seven times the test mean of 7.7\,mm, because a few missed millilitres far away
barely affect overlap but dominate boundary distance. Over all
test cases, Dice correlates only moderately with HD95 and with absolute volume error (Spearman $-$0.74 and $-$0.73), and
HD95 with volume error weakly (0.52). Overlap, boundary distance and volume answer different questions (how much agrees,
how far the worst errors reach, how wrong a volume measurement would be), so a report needs all three.

\textbf{Bonus: missing modalities.} Without retraining, replacing FLAIR by zeros collapses the model (mean test Dice
0.087; 48 of 73 predictions empty), whereas zeroing T2, T1 or T1ce costs 0.096, 0.062 and 0.028 Dice
(Appendix Table~\ref{tab:bonus}). Whole tumour is mostly oedema, which FLAIR shows best, and the network relies on it
almost entirely. Zeroing is an imperfect stand-in for a missing acquisition: after z-scoring, 0 is the brain's mean, so
the channel becomes a flat image the network never saw, whereas a real missing scan is absent, replaced by another
sequence or degraded. Training with random modality dropout and an availability mask, or with per-modality encoders
fused by an order-free operation, would let the network cope with any subset.

\textbf{Bonus: sliding-window seams.} On validation, Dice rises from 0.850 (0\% overlap, about 26 windows per case) to
0.862 (25\%) and 0.873 (50\%, about 77 windows), while the run time grows from 0.33 to 0.88\,s per case. When the 0\% and 50\% predictions are
compared, labels change for 10.9\% of compared voxels within four voxels of a window border but for
1.7\% sixteen or more voxels away (Appendix Figure~\ref{fig:seams}): border voxels see the tumour cut off by
the patch edge, and overlap lets a window in which they lie near the centre dominate. Choosing the overlap by repeatedly
checking test Dice would turn the test set into a second validation set and bias the reported score upwards; we fixed
50\% on validation.

\section{Discussion and limitations}
The pipeline reaches a whole-tumour test Dice of 0.885 with a lightweight network and a short budget, and the
controlled comparison over three seeds shows that, at this patch size and tumour size, tumour-aware sampling shifts
errors from misses to false positives without improving overlap. Five limitations bound these conclusions.
\emph{Small training budget:} each run saw 12,000 patches on one laptop GPU and validation Dice was still rising, so
absolute numbers understate what the architecture can reach; a longer schedule and a larger patch would test whether
the sampler ranking holds. \emph{No external-site validation:} all cases come from one curated benchmark with shared
skull stripping, co-registration and resampling; raw scans from another hospital or protocol may perform worse, and
only a held-out site would tell. \emph{Reference, not truth:} each case has one research segmentation, and raters
disagree most at the oedema edge, exactly where most of our errors lie. \emph{Three seeds, one split:} confidence
intervals cover case sampling, not re-splitting. \emph{Whole tumour only, and a Dice-tuned rule:} the clinically
important sub-regions were not modelled, and choosing postprocessing on Dice alone removed small true lesions; a rule
chosen jointly on Dice and lesion-wise detection would be safer.

\section{Conclusion}
A patch-trained 3D U-Net with careful geometry handling, whole-volume sliding-window inference and validation-only
choices segments the whole tumour with a test Dice of 0.885 and a median HD95 of 4.9\,mm. Tumour-aware sampling did not
help here, and the error analysis shows why single metrics mislead: the remaining errors are boundary rims, missed small
lesions and context failures, which Dice, HD95 and volume error each see differently.

\medskip\noindent\textbf{Acknowledgements and AI use.} Data: MSD Task01 \cite{simpson2019,antonelli2022}, from BraTS \cite{menze2015,bakas2017}. Software: PyTorch \cite{paszke2019}, NiBabel \cite{nibabel}; figures use the \emph{paper-figure} style code \cite{paperfigure}. AI use is disclosed in \texttt{ai\_use\_declaration.txt} in the repository.

\clearpage
\begin{thebibliography}{99}\small
\bibitem{antonelli2022} M. Antonelli et al. The Medical Segmentation Decathlon. \emph{Nature Communications}, 13:4128, 2022. doi:10.1038/s41467-022-30695-9
\bibitem{simpson2019} A. L. Simpson et al. A large annotated medical image dataset for the development and evaluation of segmentation algorithms. arXiv:1902.09063, 2019.
\bibitem{menze2015} B. H. Menze et al. The Multimodal Brain Tumor Image Segmentation Benchmark (BRATS). \emph{IEEE Transactions on Medical Imaging}, 34(10):1993--2024, 2015. doi:10.1109/TMI.2014.2377694
\bibitem{bakas2017} S. Bakas, H. Akbari, A. Sotiras, M. Bilello, M. Rozycki, J. S. Kirby, J. B. Freymann, K. Farahani, C. Davatzikos. Advancing The Cancer Genome Atlas glioma MRI collections with expert segmentation labels and radiomic features. \emph{Scientific Data}, 4:170117, 2017. doi:10.1038/sdata.2017.117
\bibitem{ronneberger2015} O. Ronneberger, P. Fischer, T. Brox. U-Net: Convolutional networks for biomedical image segmentation. In \emph{MICCAI 2015}, LNCS 9351, pages 234--241. Springer, 2015. doi:10.1007/978-3-319-24574-4\_28
\bibitem{cicek2016} \"O. \c{C}i\c{c}ek, A. Abdulkadir, S. S. Lienkamp, T. Brox, O. Ronneberger. 3D U-Net: Learning dense volumetric segmentation from sparse annotation. In \emph{MICCAI 2016}, LNCS 9901, pages 424--432. Springer, 2016. doi:10.1007/978-3-319-46723-8\_49
\bibitem{ulyanov2016} D. Ulyanov, A. Vedaldi, V. Lempitsky. Instance normalization: The missing ingredient for fast stylization. arXiv:1607.08022, 2016.
\bibitem{milletari2016} F. Milletari, N. Navab, S.-A. Ahmadi. V-Net: Fully convolutional neural networks for volumetric medical image segmentation. In \emph{2016 Fourth International Conference on 3D Vision (3DV)}, pages 565--571. IEEE, 2016. doi:10.1109/3DV.2016.79
\bibitem{loshchilov2019} I. Loshchilov, F. Hutter. Decoupled weight decay regularization. In \emph{ICLR}, 2019. arXiv:1711.05101
\bibitem{isensee2021} F. Isensee, P. F. Jaeger, S. A. A. Kohl, J. Petersen, K. H. Maier-Hein. nnU-Net: a self-configuring method for deep learning-based biomedical image segmentation. \emph{Nature Methods}, 18(2):203--211, 2021. doi:10.1038/s41592-020-01008-z
\bibitem{taha2015} A. A. Taha, A. Hanbury. Metrics for evaluating 3D medical image segmentation: analysis, selection, and tool. \emph{BMC Medical Imaging}, 15:29, 2015. doi:10.1186/s12880-015-0068-x
\bibitem{maierhein2024} L. Maier-Hein, A. Reinke, et al. Metrics reloaded: recommendations for image analysis validation. \emph{Nature Methods}, 21(2):195--212, 2024. doi:10.1038/s41592-023-02151-z
\bibitem{paszke2019} A. Paszke et al. PyTorch: An imperative style, high-performance deep learning library. In \emph{Advances in Neural Information Processing Systems 32}, 2019.
\bibitem{nibabel} C. J. Markiewicz, M. Brett, et al. NiBabel (software; version 5.4.2 used, released 2026). All-versions Zenodo record, doi:10.5281/zenodo.591597
\bibitem{paperfigure} \emph{paper-figure} figure-style skill (supplied as a Claude Code skill); its Appendix A plotting code is reused in \texttt{helpers/paper\_plot\_style.py}.
\end{thebibliography}


\clearpage
\appendix
\section*{Appendix}
\noindent\begin{minipage}{\textwidth}
\centering\fig{C_augmentation.pdf}
\captionof{figure}{\textbf{Augmentation changes position and intensity but keeps image and label aligned.} One training
patch (rule: first uniform patch with both a flip and a rotation whose middle slice shows tumour and the brain's edge)
without (a, b) and with (d, e) augmentation; gold: reference whole tumour, transformed with nearest-neighbour
interpolation. (c) FLAIR intensities of the brain voxels in the patch.}
\label{fig:augment}
\end{minipage}

\bigskip\noindent\begin{minipage}{\textwidth}
\centering\fig{F_volumes.pdf}
\captionof{figure}{\textbf{Predicted volumes match the reference without bias but with a wide spread for single cases.}
Test split, 73 cases, final model after postprocessing. (a) Predicted against reference volume. (b) Bland--Altman plot:
volume error as a percentage of the reference volume; solid line: mean (bias), dashed: bias $\pm$ 1.96 SD. Rings: the
cases of Figure~\ref{fig:errors}.}
\label{fig:volumes}
\end{minipage}

\bigskip\noindent\begin{minipage}{\textwidth}
\centering\footnotesize
\captionof{table}{\textbf{Bonus tests.} Top: test split, 73 cases, one channel replaced by zeros at test time (final
model, postprocessed); $\Delta$Dice paired with the four-channel result. Bottom: window overlap, final model,
postprocessed; s/case: inference time per case on an Apple M4 Max GPU.}
\label{tab:bonus}
\begin{tabular}{lccccc}
\toprule
input & Dice (mean) & HD95 median (mm) & $|$vol.\ error$|$ median (\%) & empty & $\Delta$Dice [95\% CI] \\
\midrule
all four channels & 0.885 & 4.9 & 10.1 & 0 & -- \\
FLAIR zeroed & 0.087 & 32.2 & 100.0 & 48 & $-$0.798 [$-$0.830, $-$0.762] \\
T1 zeroed & 0.823 & 14.5 & 19.6 & 0 & $-$0.062 [$-$0.084, $-$0.042] \\
T1ce zeroed & 0.857 & 5.8 & 17.1 & 0 & $-$0.028 [$-$0.041, $-$0.018] \\
T2 zeroed & 0.789 & 7.5 & 19.4 & 2 & $-$0.096 [$-$0.140, $-$0.057] \\
\bottomrule
\end{tabular}\\[4pt]
\begin{tabular}{llccccc}
\toprule
split & overlap & windows/case & s/case & Dice (mean) & HD95 median (mm) & $|$vol.\ error$|$ median (\%) \\
\midrule
test & 0\% & 26.9 & 0.32 & 0.868 & 4.7 & 12.0 \\
test & 25\% & 34.9 & 0.39 & 0.881 & 5.0 & 10.8 \\
test & 50\% & 77.8 & 0.90 & 0.885 & 4.9 & 10.1 \\
validation & 0\% & 26.4 & 0.33 & 0.850 & 5.9 & 12.7 \\
validation & 25\% & 34.9 & 0.39 & 0.862 & 5.9 & 11.6 \\
validation & 50\% & 76.7 & 0.88 & 0.873 & 5.1 & 10.2 \\
\bottomrule
\end{tabular}
\end{minipage}

\bigskip\noindent\begin{minipage}{\textwidth}
\centering\footnotesize
\captionof{table}{\textbf{The three splits have similar tumour volumes.} Reference whole-tumour volume per split (voxel
count $\times$ 1\,mm$^3$); IQR: interquartile range. Test labels were read only to count voxels.}
\label{tab:data}
\begin{tabular}{lrccc}
\toprule
split & cases & median [IQR] (mL) & mean $\pm$ SD (mL) & range (mL) \\
\midrule
train & 338 & 95.5 [60.0, 143.9] & 103.0 $\pm$ 57.0 & 8.5--265.3 \\
validation & 73 & 91.6 [51.9, 148.1] & 100.8 $\pm$ 63.0 & 7.3--266.9 \\
test & 73 & 89.1 [54.4, 157.7] & 107.4 $\pm$ 66.3 & 10.1--318.3 \\
\bottomrule
\end{tabular}\\[6pt]
\fig{B_volumes.pdf}
\captionof{figure}{\textbf{Tumour volume distributions overlap across splits.} Dots: cases; black bar: interquartile
range; red tick: median.}
\label{fig:split}
\end{minipage}

\bigskip\noindent\begin{minipage}{\textwidth}
\centering\fig{X2_seams.pdf}
\captionof{figure}{\textbf{Overlap helps most at window borders.} (a) Validation Dice against inference time for 0, 25 and
50\% overlap. (b) Share of voxels whose label differs between 0\% and 50\% overlap, by distance to the nearest window
border of the 0\% grid (validation; voxels where either map has $p \ge 0.05$). (c) $|\Delta p|$ between the two maps
on one validation case, with the 0\% window borders (dashed).}
\label{fig:seams}
\end{minipage}

\end{document}
